\documentclass[11pt,aspectratio=169]{beamer}

\usepackage{iftex}
\ifPDFTeX
\usepackage[utf8]{inputenc}
\usepackage[T1]{fontenc}
\fi
\usepackage[slovak]{babel}
\usepackage{amsmath}
\usepackage{amssymb}
    \usepackage{graphicx}
\usepackage{tikz}
\usepackage{qtree}
\usepackage{forest}

\usetikzlibrary{arrows.meta,positioning}

\usetheme[numbering=fraction]{metropolis}
\setbeamertemplate{navigation symbols}{}
\newcommand{\slidefootnote}{\insertshortauthor\hspace{0.6em}|\hspace{0.6em}\insertshorttitle}
\setbeamertemplate{frame footer}{\slidefootnote}

\newcommand{\UNION}{\textbf{UNION}} \newcommand{\JOIN}{\textbf{JOIN}}
\newcommand{\ANTIJOIN}{\textbf{ANTIJOIN}} \newcommand{\REC}{\textbf{REC}}
\newcommand{\DEF}{\textbf{DEF}} \newcommand{\tp}{\textbf{t}\space}

\title[Kompilátor Datalogu do relačnej algebry]{Kompilátor Datalogu s funkčnými symbolmi do relačnej algebry}
\author[Matúš Duchyňa]{Matúš Duchyňa}
\institute[]{%
    Fakulta matematiky, fyziky a informatiky\\
    Univerzita Komenského v Bratislave\\[0.4em]
    Vedúci práce: doc. Mgr. Tomáš Plachetka, Dr.
}
\date{8. 6., 2026}

\newcommand{\antijoin}{\mathbin{\triangleright}}

\begin{document}

\begin{frame}
    \titlepage
\end{frame}

\begin{frame}{Motivácia}
    \begin{itemize}
        \item Experimentálny databázový systém s relačnou algebrou ako dotazovacím jazykom
        \item Relačná algebra
        \begin{itemize}
            \item Umožnuje určiť vykonávací plán
            \item Nízkoúrovňová
        \end{itemize}
        \item Datalog ako nový dotazovací jazyk
    \end{itemize}
\end{frame}

\begin{frame}{Ciele práce}
    \begin{block}{Hlavný cieľ}
        Kompilátor z Datalogu do relačnej algebry
    \end{block}
    \begin{itemize}
        \item Databázový systém môže obsahovať tabuľky
        \begin{itemize}
            \item Vymenovaním prvkov
            \item Popísané nejakým programom
        \end{itemize}
        \item Kompilátor ako samostatný nástroj, bez závislosti na konkrétnom databázovom systéme
        \item Podpora funkčných symbolov
    \end{itemize}
\end{frame}

\begin{frame}{Datalog}
    \begin{itemize}
        \item Tabuľky - predikáty
        \item Vstavané tabuľky - vstavané predikáty (označené s \$ pred názvom)
    \end{itemize}
    
    Pod pojmom \textbf{parameter} rozumieme:
    \begin{enumerate}
        \item ľubovoľný term
        \item ľubovoľný predikát alebo vstavaný predikát
        \item ľubovoľný zoznam parametrov, ktorý zapisujeme ako $[P_1, P_2, \dots, P_n]$,\\ kde
        $P_1, P_2, \dots, P_n$ sú parametre
    \end{enumerate}
\end{frame}

\begin{frame}{Relačná algebra}
    \begin{itemize}
        \item Tabuľky - relácie
        \item Vstavané \textit{jednoduché} tabuľky - vstavané relácie (označené s \$ pred názvom)
        \item Vstavané \textit{zložité} tabuľky - vstavané operátory (označené s \$ pred názvom)
    \end{itemize}
    
    \begin{columns}[T,onlytextwidth]
        \begin{column}{0.48\textwidth}
            \textbf{Klasická relačná algebra}
            \begin{itemize}
                \item Množinové operátory, $\times$
                \item $\Pi, \sigma, \rho$
                \item $\Join, \antijoin, \Join_{C}$
            \end{itemize}
        \end{column}
        \begin{column}{0.48\textwidth}
            \textbf{Naša relačná algebra}
            \begin{itemize}
                \item \UNION
                \item \JOIN
                \item \ANTIJOIN
                \item \REC
            \end{itemize}
        \end{column}
    \end{columns}

\end{frame}

\begin{frame}{Operátory}
    \textbf{Parameter} (\textit{v kontexte relačnej algebry}) je:
    \begin{enumerate}
        \item ľubovoľný term
        \item ľubovoľnú \textit{reláciu} alebo vstavaný operátor (ktorého výstupom je relácia)
        \item ľubovoľný zoznam parametrov, ktorý zapisujeme ako $[P_1, P_2, \dots, P_n]$,\\ kde
        $P_1, P_2, \dots, P_n$ sú parametre
    \end{enumerate}
    
    \medskip
    
    \textbf{Operátor} je $\$O(P_1, P_2, \dots, P_n)$, kde $\$O$ je názov operátora a\\
    $P_1, P_2, \dots, P_n$ sú parametre operátora. Výstupom operátora je relácia.
\end{frame}

\begin{frame}{Syntaktická analýza}
    \begin{columns}[T,onlytextwidth]
        \begin{column}{0.25\textwidth}
            \scriptsize
            \begin{align*}
                G &= (N, T, P, \sigma)\text{, kde}\\
                N &= \{\sigma, \alpha\}\text{,}\\
                T &= \{\texttt{+}, \cdot, \texttt{(}, \texttt{)}, 0, \dots, 9\}\text{,}\\
                P &= \{\sigma \rightarrow \texttt{(}\sigma\texttt{)} \, | \, \sigma \texttt{+} \sigma \, | \, \sigma \cdot \sigma\\
                  &\quad \quad\, | \, 1 \alpha \, | \dots | \, 9 \alpha \, | \, 0,\\
                  &\quad \quad \alpha \rightarrow 0 \alpha \, | \dots | \, 9 \alpha \, | \, \varepsilon\}
            \end{align*}
            \normalsize
        \end{column}
        \begin{column}{0.75\textwidth}
            \begin{center}
            \scriptsize
            Odvodenie slova \texttt{(4 + 5) $\cdot$ 123} $\cdot$ 0
            \[
            \Tree
            [.$\sigma$ 
                [.$\sigma$ 
                    $\texttt{(}$
                    [.$\sigma$
                        [.$\sigma$
                            $4$
                            [.$\alpha$
                                $\varepsilon$
                            ]
                        ]
                        $\texttt{+}$
                        [.$\sigma$
                            $5$
                            [.$\alpha$
                                $\varepsilon$
                            ]
                        ]
                    ]
                    $\texttt{)}$
                ] 
                $\cdot$ 
                [.$\sigma$ 
                    [.$\sigma$
                        $1$
                        [.$\alpha$
                            $2$
                            [.$\alpha$
                                $3$
                                [.$\alpha$
                                    $\varepsilon$
                                ]
                            ]
                        ]
                    ]
                    $\cdot$
                    [.$\sigma$
                        $0$
                    ]
                ]
            ]
            \]
            \normalsize
            \end{center}
        \end{column}
    \end{columns}
\end{frame}

\begin{frame}{Tokeny}
    \scriptsize
    $$\texttt{INT} = \{0\} \cup \{1, 2, 3, 4, 5, 6, 7, 8, 9\} \cdot \{1, 2, 3, 4, 5, 6, 7, 8, 9,
    0\}^* \approx \mathbb{N}_0$$
    $$\texttt{LPAREN}(\texttt{(}),\, \texttt{INT}(4),\, \texttt{PLUS}(\texttt{+}),\, \texttt{INT}(5),\,
\texttt{RPAREN}(\texttt{)}),\, \texttt{MULT}(\cdot),\, \texttt{INT}(123),\, \texttt{MULT}(\cdot),\,
\texttt{INT}(0)$$

\begin{columns}[T,onlytextwidth]
        \begin{column}{0.50\textwidth}
            \begin{align*}
                    G' &= (N', T', P', \beta)\text{, kde}\\
                    N' &= \{\beta\}\text{,}\\
                    T' &= \{\texttt{+}, \cdot, \texttt{(}, \texttt{)}, \mathtt{INT}\}\text{,}\\
                    P' &= \{\beta \rightarrow \texttt{(}\beta\texttt{)} \, | \, \beta \texttt{+} \beta \, | \, \beta \cdot \beta
                                    \, | \, \texttt{INT}\}
            \end{align*}
        \end{column}
        \begin{column}{0.50\textwidth}
            \begin{center}
                \[
                \Tree
                [.$\beta$ 
                    [.$\beta$ 
                        $\texttt{(}$
                        [.$\beta$
                            [.$\beta$
                                \qroof{$4$}.$\texttt{INT}$ 
                            ]
                            $\texttt{+}$
                            [.$\beta$
                                \qroof{$5$}.$\texttt{INT}$ 
                            ]
                        ]
                        $\texttt{)}$
                    ] 
                    $\cdot$ 
                    [.$\beta$ 
                        [.$\beta$
                            \qroof{$123$}.$\texttt{INT}$ 
                        ]
                        $\cdot$
                        [.$\beta$
                            \qroof{$0$}.$\texttt{INT}$ 
                        ]
                    ]
                ]
                \]
                \normalsize
            \end{center}
        \end{column}
    \end{columns}
\end{frame}

\begin{frame}{Algoritmus prekladu}
    \begin{enumerate}
    \item Vytvorenie lexeru a token streamu
    \item Vytvorenie parseru a stromu odvodenia
    \item Generácia AST z odvodenia
    \item Preklad predikátov do relačnej algebry s nevyplnenými referenciami
    \begin{enumerate}
        \item Preklad definícií pomocou \JOIN\ a \UNION\
        \item Preklad klauzúl pomocou \JOIN\ a \ANTIJOIN\
        \item Spojenie klauzúl do jednej relácie pomocou \UNION\
    \end{enumerate}
    \item Zostavenie výsledného výrazu v relačnej algebre
    \begin{enumerate}
        \item Zostavenie grafu závislosti predikátov
        \item Zabalenie rekurzívnych predikátov do \REC\
        \item Nahradenie referencií výrazmi, ktoré im prislúchajú
        \item Zabalenie výsledného výrazu do \JOIN-u reprezentujúceho dotaz
    \end{enumerate}
\end{enumerate}
\end{frame}

\begin{frame}{1. a 2. krok}
    \begin{center}
        ANTLR
    \end{center}
    
    \begin{columns}[T,onlytextwidth]
        \begin{column}{0.48\textwidth}
            \footnotesize
            \texttt{grammar datalog;\\
                    QUERY\_MARKER: '?-';\\
                    BODY\_HEAD\_SEP: ':-';\\
                    NOT: '\textbackslash+';\\
                    BUILT\_IN: '\$';\\
                    program: (clause | definition)* query;\\
                    query: QUERY\_MARKER predicate EOL;\\
                    clause: name '(' ')' BODY\_HEAD\_SEP\\
                    \hspace{2em} predicate (',' predicate)* EOL\\
                    \hspace{1em} | name '(' term (',' term)* ')' \\
                    \hspace{2em} BODY\_HEAD\_SEP predicate\\
                    \hspace{2em} (',' predicate)* EOL;
            }
        \end{column}
        \begin{column}{0.48\textwidth}
            \footnotesize
            \texttt{definition: name BODY\_HEAD\_SEP\\
                \hspace{2em} '\{' term\_tuple (',' term\_tuple)* '\}' EOL;\\
                term: variable\\
                \hspace{1em} | function '(' term (',' term)* ')';\\
                name: LOWER (LOWER | NUMBER | '\_')*;\\
                function: LOWER (LOWER | NUMBER | '\_')*;\\
                variable: UPPER (UPPER | NUMBER | '\_')*;\\
                UPPER: [A-Z]+;\\
                LOWER: [a-z]+;\\
                NUMBER: [0-9]+;\\
                WS: [ \textbackslash t\textbackslash r\textbackslash n]+ -\textgreater\ skip\\
            }
        \end{column}
    \end{columns}
\end{frame}

\begin{frame}{3. Generácia AST z odvodenia}
\scriptsize
\begin{columns}[T,onlytextwidth]
    \begin{column}{0.25\textwidth}
        \begin{center}
        \[
        \Tree
        [.\texttt{predicate}
            [.\texttt{normal\_predicate}
                \texttt{name}
                \texttt{(}
                \textit{possible terms}
                \texttt{)}
            ]
        ]
        \]
        \textit{Obyčajný predikát}
        \end{center}
    \end{column}
    \begin{column}{0.60\textwidth}
        \begin{center}
        \[
        \Tree
        [.\texttt{built\_in\_predicate}
            [.\texttt{negative\_built\_in\_predicate}
                \qroof{\texttt{\textbackslash+}}.\texttt{NOT}
                [.\texttt{normal\_built\_in\_predicate}
                    \qroof{\texttt{\$}}.\texttt{BUILT\_IN}
                    \texttt{name}
                    \texttt{(}
                    \textit{possible parameters}
                    \texttt{)}
                ]
            ]
        ]
        \]
        \textit{Negovaný vstavaný predikát}
        \end{center}
    \end{column}
\end{columns}
\end{frame}

\begin{frame}{AST}
\centering
\resizebox{!}{0.85\textheight}{%
\begin{forest}
  for tree={
    draw,
    align=left,
    font=\ttfamily\footnotesize,
    inner sep=5pt,
    fill=gray!10,
    edge={-{Stealth[scale=0.8]}},
    l sep=1.2cm,
    s sep=0.3cm,
  }
  [{\bfseries DatalogProgram\\ - clauses: List<Clause>\\ - definitions: List<Definition>\\ - query: Query}
  [{\bfseries Clause\\ - head: Predicate\\ - body: List<Atom>}
    [{\bfseries Predicate\\ - name: String\\ - terms: List<Term>\\ - isBuiltIn: bool}
        [{\bfseries (Term) Constant\\ - value: String},       edge={dashed,-{Triangle[open,scale=0.8]}}]
        [{\bfseries (Term) Variable\\ - name: String},        edge={dashed,-{Triangle[open,scale=0.8]}}]
        [{\bfseries (Term) FunctionTerm\\ - name: String\\ - args: List<Term>}, edge={dashed,-{Triangle[open,scale=0.8]}}]
        [{\bfseries (Term) Parameter\\ - text: String},       edge={dashed,-{Triangle[open,scale=0.8]}}]
    ]
    [{\bfseries Atom\\ - predicate: Predicate\\ - isNegated: boolean}]
  ]
  [{\bfseries Definition\\ - name: String\\ - arity: int\\ - tuples: List<List<Term>>}]
  [{\bfseries Query\\ - predicate: Predicate}]
  ]
\end{forest}%
}
\end{frame}

\begin{frame}{4. Preklad predikátov do relačnej algebry s nevyplnenými referenciami}
\scriptsize
\begin{columns}[T,onlytextwidth]
    \begin{column}{0.32\textwidth}
        \textbf{1) Definície}\\
        \texttt{color :- \{(r()), (b())\}.}
        
        \vspace{2,5em}
        sa preloží ako
        \vspace{1em}
        
        \texttt{UNION([\\
            \hspace{1em} JOIN([\$true], [[]], [r()]),\\
            \hspace{1em} JOIN([\$true], [[]], [b()])\\
            ])
        }
    \end{column}
    \begin{column}{0.32\textwidth}
        \textbf{2) Klauzuly}\\
        \texttt{t(X, Z) :- e(X, Y), e(Y, Z),\\
            \hspace{4em} \textbackslash+fn(Y), \textbackslash+fe(X, Y).
        }
        
        \vspace{1em}
        sa preloží ako
        \vspace{1em}
        
        \texttt{ANTIJOIN([\\
            \hspace{1em} JOIN([e, e], \\
            \hspace{2em} [[X, Y], [Y, Z]], \\
            \hspace{2em} [X, Y, Z]\\
            \hspace{1em} ),\\
            \hspace{1em} fn, fe],\\
            \hspace{1em} [[X, Y, Z], [Y], [X, Y]],\\
            \hspace{1em} [X, Z]\\
            ])
        }
    \end{column}
    \begin{column}{0.32\textwidth}
        \textbf{3) Zjednotenie výsledkov}\\
        \texttt{p(X) :- \$a(X).\\
            p(Y) :- \$b(Y).
        }
        
        \vspace{1em}
        sa preloží ako
        \vspace{1em}
        
        \texttt{UNION([\\
            \hspace{1em} JOIN([\$a], [[X]], [X]),\\
            \hspace{1em} JOIN([\$b], [[Y]], [Y])\\
            ])
        }
    \end{column}
\end{columns}
\end{frame}

\begin{frame}{Poradie atómov v klauzule}
    2 sémanticky ekvivalentné programy:
    \vspace{1em}
    
    \begin{columns}[T,onlytextwidth]
        \begin{column}{0.49\textwidth}
            \texttt{p(Z) :- \$do10(X), \$do10(Y),\\
                \hspace{3.5em} \$plus(X, Y, Z).
            }
            
            \vspace{3em}
            
            \texttt{JOIN(\\
                \hspace{1em} [\$do10, \$do10, \$plus],\\
                \hspace{1em} [[X], [Y], [X, Y, Z]],\\
                \hspace{1em} [Z]\\
                )
            }
        \end{column}
        \begin{column}{0.49\textwidth}
            \texttt{p(Z) :- \$plus(X, Y, Z),\\
                \hspace{3.5em} \$do10(X), \$do10(Y).
            }
            
            \vspace{3em}
            
            \texttt{JOIN(\\
                \hspace{1em} [\$plus, \$do10, \$do10],\\
                \hspace{1em} [[X, Y, Z], [X], [Y]],\\
                \hspace{1em} [Z]\\
                )
            }
        \end{column}
    \end{columns}
\end{frame}

\begin{frame}{5. Zostavenie výsledného výrazu v relačnej algebre}
    \begin{columns}[T,onlytextwidth]
        \begin{column}{0.68\textwidth}
            \begin{itemize}
                \item \texttt{p -> JOIN([q, r], [[X], [X]], [X])}
                \item \texttt{q -> JOIN([p, s], [[X], [X]], [X])}
                \item \texttt{r -> JOIN([t, t], [[X], [X]], [X])}
                \item \texttt{s -> UNION([...])}
                \item \texttt{t -> UNION([...])}
            \end{itemize}
        \end{column}
        \begin{column}{0.28\textwidth}
            \begin{center}
                \begin{tikzpicture}[
                    node distance=2cm,
                    every node/.style={draw, align=center, fill=gray!10, font=\ttfamily\footnotesize},
                    edge/.style={-{Stealth[scale=0.8]}, thick},
                ]
                    \node (p) {p};
                    \node (q) [below left of=p] {q};
                    \node (r) [below right of=p] {r};
                    \node (s) [below of=q] {s};
                    \node (t) [below of=r] {t};
                    
                    \draw[edge] (p) -- (q);
                    \draw[edge] (q) -- (p);
                    \draw[edge] (p) -- (r);
                    \draw[edge] (q) -- (s);
                    \draw[edge] (r) -- (t);
                \end{tikzpicture}
            \end{center}
        \end{column}
    \end{columns}
\end{frame}

\begin{frame}{5. Zostavenie výsledného výrazu v relačnej algebre}
    \begin{columns}[T,onlytextwidth]
        \begin{column}{0.68\textwidth}
            \begin{itemize}
                \item \texttt{p -> REC(p, JOIN([q, r], [[X], [X]], [X]))}
                \item \texttt{q -> REC(q, JOIN([p, s], [[X], [X]], [X]))}
                \item \texttt{r -> JOIN([t, t], [[X], [X]], [X])}
                \item \texttt{s -> UNION([...])}
                \item \texttt{t -> UNION([...])}
            \end{itemize}
        \end{column}
        \begin{column}{0.28\textwidth}
            \begin{center}
                \begin{tikzpicture}[
                    node distance=2cm,
                    every node/.style={draw, align=center, fill=gray!10, font=\ttfamily\footnotesize},
                    edge/.style={-{Stealth[scale=0.8]}, thick},
                ]
                    \node (p) {p};
                    \node (q) [below left of=p] {q};
                    \node (r) [below right of=p] {r};
                    \node (s) [below of=q] {s};
                    \node (t) [below of=r] {t};
                    
                    \draw[edge] (p) -- (q);
                    \draw[edge] (q) -- (p);
                    \draw[edge] (p) -- (r);
                    \draw[edge] (q) -- (s);
                    \draw[edge] (r) -- (t);
                \end{tikzpicture}
            \end{center}
        \end{column}
    \end{columns}
\end{frame}

% 1
\begin{frame}{Zostavenie dotazu \texttt{?- q(f(X)).}}
    \texttt{REC(q, JOIN([\\
        \hspace{1em} p,\\
        \hspace{1em} s\\
        ], [[X], [X]], [X]))
    }
\end{frame}

% 2
\begin{frame}{Zostavenie dotazu \texttt{?- q(f(X)).}}
    \texttt{REC(q, JOIN([\\
        \hspace{1em} p,\\
        \hspace{1em} UNION([...])\\
        ], [[X], [X]], [X]))
    }
\end{frame}

% 3
\begin{frame}{Zostavenie dotazu \texttt{?- q(f(X)).}}
    \texttt{REC(q, JOIN([\\
        \hspace{1em} REC(p, JOIN([\\
        \hspace{2em} q,\\
        \hspace{2em} r\\
        \hspace{1em} ], [[X], [X]], [X])),\\
        \hspace{1em} UNION([...])\\
        ], [[X], [X]], [X]))
    }
\end{frame}

% 4
\begin{frame}{Zostavenie dotazu \texttt{?- q(f(X)).}}
    \texttt{REC(q, JOIN([\\
        \hspace{1em} REC(p, JOIN([\\
        \hspace{2em} q,\\
        \hspace{2em} JOIN([\\
        \hspace{3em} t,\\
        \hspace{3em} t\\
        \hspace{2em} ]], [[X], [X]], [X])\\
        \hspace{1em} ], [[X], [X]], [X])),\\
        \hspace{1em} UNION([...])\\
        ], [[X], [X]], [X]))
    }
\end{frame}

% 5
\begin{frame}{Zostavenie dotazu \texttt{?- q(f(X)).}}
    \texttt{REC(q, JOIN([\\
        \hspace{1em} REC(p, JOIN([\\
        \hspace{2em} q,\\
        \hspace{2em} JOIN([\\
        \hspace{3em} UNION([...]),\\
        \hspace{3em} UNION([...])\\
        \hspace{2em} ]], [[X], [X]], [X])\\
        \hspace{1em} ], [[X], [X]], [X])),\\
        \hspace{1em} UNION([...])\\
        ], [[X], [X]], [X]))
    }
\end{frame}

% 6
\begin{frame}{Zostavenie dotazu \texttt{?- q(f(X)).}}
    \texttt{JOIN([\\
        \hspace{1em} REC(q, JOIN([\\
        \hspace{2em} REC(p, JOIN([\\
        \hspace{3em} q,\\
        \hspace{3em} JOIN([\\
        \hspace{4em} UNION([...]),\\
        \hspace{4em} UNION([...])\\
        \hspace{3em} ]], [[X], [X]], [X])\\
        \hspace{2em} ], [[X], [X]], [X])),\\
        \hspace{2em} UNION([...])\\
        \hspace{1em} ], [[X], [X]], [X]))\\
        ], [[f(X)]], [f(X)])
    }
\end{frame}

\begin{frame}{}
    \begin{center}
        Ďakujem za pozornosť!
    \end{center}
\end{frame}

\end{document}